\paragraph{Simulation.} Multi-excitation fluorescence hyperspectral images were simulated with SpectralForge Studio (version 0.0.0-test, commit 0123456789), a physics-based forward model that combines ray tracing of the scene, Monte Carlo light transport, dye absorption, fluorescence emission and a detector model. The experiment 'Tissue under collagen (wedge)' (SHA-256 \texttt{c9e84764a8fc81c7}) was simulated in sample mode.

\paragraph{Scene.} The scene consisted of 1 object(s) on a 600 mm background plane of black\_anodized. Tissue / collagen block: a layered block of 40 × 40 mm footprint with layers (bottom to top) substrate (tissue\_substrate, 5\,mm), tissue (tissue\_fluor\_layer, 1\,mm) then collagen (collagen\_cover, 1.5\,mm to 0.05\,mm along x), centred at (0.0, 0.0, 0.0) mm.

\paragraph{Optical properties.} Optical properties at the excitation wavelengths (310\,nm, 325\,nm, 340\,nm, 365\,nm, 385\,nm, 400\,nm, 415\,nm and 430\,nm) are listed below: refractive index n, scattering anisotropy g, absorption coefficient $\mu_a$ and reduced scattering coefficient $\mu_s'$ = $\mu_s(1-g)$ in mm$^{-1}$, and dye concentrations. Non-scattering (clear) materials have no g or $\mu_s'$.

\begin{tabular}{lrrrrrrrrrrl}
\hline material & n & g & $\mu_a$ / $\mu_s'$ at 310 nm (mm$^{-1}$) & $\mu_a$ / $\mu_s'$ at 325 nm (mm$^{-1}$) & $\mu_a$ / $\mu_s'$ at 340 nm (mm$^{-1}$) & $\mu_a$ / $\mu_s'$ at 365 nm (mm$^{-1}$) & $\mu_a$ / $\mu_s'$ at 385 nm (mm$^{-1}$) & $\mu_a$ / $\mu_s'$ at 400 nm (mm$^{-1}$) & $\mu_a$ / $\mu_s'$ at 415 nm (mm$^{-1}$) & $\mu_a$ / $\mu_s'$ at 430 nm (mm$^{-1}$) & dyes \\ \hline
tissue\_substrate & 1.38 & 0.85 & 0.32 / 2.66 & 0.394 / 2.52 & 0.436 / 2.38 & 0.37 / 2.19 & 0.423 / 2.05 & 0.648 / 1.96 & 1.23 / 1.88 & 0.841 / 1.8 & NADH 8 \textmu{}M and collagen 80 \textmu{}M \\
tissue\_fluor\_layer & 1.38 & 0.85 & 0.14 / 2.13 & 0.192 / 2.01 & 0.22 / 1.91 & 0.167 / 1.75 & 0.173 / 1.64 & 0.274 / 1.57 & 0.543 / 1.5 & 0.444 / 1.44 & NADH 120 \textmu{}M and FAD 60 \textmu{}M \\
collagen\_cover & 1.41 & 0.9 & 1.17 / 3.55 & 2.01 / 3.35 & 2.45 / 3.18 & 1.63 / 2.92 & 0.62 / 2.74 & 0.209 / 2.61 & 0.0521 / 2.5 & 0.00969 / 2.4 & collagen 600 \textmu{}M and elastin 60 \textmu{}M \\
black\_anodized & 1.5 & 0.6 & 64.5 / 2.54 & 61.5 / 2.48 & 58.8 / 2.43 & 54.8 / 2.34 & 51.9 / 2.28 & 50 / 2.24 & 48.2 / 2.2 & 46.5 / 2.16 & none \\
\hline \end{tabular}

\paragraph{Source.} Excitation was provided by Tunable light source (Xe + monochromator) (tunable lamp; line width 10 nm FWHM; gaussian beam of 10 mm diameter, 7° half-angle divergence; clean-up filter OD 6, step-to-step power drift 0.8 % (1σ)) at (-200.0, 240.0, 60.0) mm aimed at (0.0, 0.0, 0.0) mm. The excitation sequence (wavelength, delivered power, exposure) was 310 nm (0.42 mW, 4000 ms), 325 nm (0.85 mW, 2000 ms), 340 nm (1.3 mW, 1000 ms), 365 nm (2.6 mW, 250 ms), 385 nm (3.4 mW, 60 ms), 400 nm (3.9 mW, 50 ms), 415 nm (4.2 mW, 80 ms) and 430 nm (4.5 mW, 80 ms).

\paragraph{Detection.} Images were recorded with Nuance-class HSI camera: a 35 mm f/4 lens focused at 330 mm, a 348×256 px sensor (6.45 µm pitch, 4×4 binning, effective 25.8 µm; ccd_bv quantum-efficiency curve, 16-bit, gain 1.5 e⁻/ADU, offset 100 ADU, read noise 8 e⁻, dark current 0.5 e⁻/s, full well 60000 e⁻, PRNU 1 %, DSNU 2 e⁻) and a lctf spectral filter (31 bands in 10 nm steps starting 20 nm above each excitation line; 20 nm FWHM, 415 nm long-pass (OD 5)).

\paragraph{Physics and Monte Carlo.} The following physics modules were switched OFF: \texttt{indirect\_bounce}; all others were enabled. The per-material transport Monte Carlo used 200,000 photons per bake. Layered blocks were solved with a 3D layered Monte Carlo (2,000,000 excitation and 2,000,000 emission photons per step, 8 sublayers per layer, lateral cell 0.5 mm, exit cell 0.5 mm, emission nodes every 20 nm, inclusion voxel 0.2 mm). Random numbers were seeded with 0 (the identical seed reproduces the identical result).

\paragraph{Study design.} The study varied cover\_mm (\texttt{/objects/0/shape/layers/2/thickness/thickness\_mm}: 0.1 to 1) and scale (\texttt{/acquisition/preview\_scale}: one of 0.2 or 0.3) using Saltelli's sampling scheme (Saltelli, 2010) with N = 64 base samples from a scrambled Sobol' sequence, i.e. N(k + 2) = 256 design points, analysed with the Saltelli (first-order) and Jansen (total-order) estimators (Jansen, 1999) with bootstrap 95 % confidence intervals (design seed 2). Each design point was simulated 3 time(s) in preview-resolution single-excitation mode with Monte Carlo seeds base+0…base+2. The recorded metrics were signal (band sum) and ratio (band ratio).

\paragraph{Software.} Simulations ran with Python 3.11.0, NumPy 2.1.3 and Numba 0.63.1 on test-platform. If you use this simulator please cite: N. Meloyan, "SpectralForge: A Physics-Based Simulator for Fluorescence Hyperspectral Imaging," 16th IEEE GRSS Workshop on Hyperspectral Imaging and Signal Processing (WHISPERS), Glasgow, 2026.
